\documentclass[11pt]{article}
\usepackage[margin=1in]{geometry}
\usepackage[T1]{fontenc}
\usepackage[utf8]{inputenc}
\usepackage{lmodern}
\usepackage{hyperref}
\usepackage{enumitem}
\usepackage{booktabs}
\usepackage{longtable}
\usepackage{xcolor}
\usepackage{listings}
\hypersetup{colorlinks=true, linkcolor=blue!60!black, urlcolor=blue!60!black}
\lstset{basicstyle=\ttfamily\small, breaklines=true, frame=single, columns=fullflexible}

\newcommand{\statusbox}[2]{
\begin{center}
\fcolorbox{blue!40!black}{blue!5}{\parbox{0.92\textwidth}{
\textbf{#1}\\[4pt]#2
}}
\end{center}
}
\newcommand{\warningbox}[2]{
\begin{center}
\fcolorbox{red!50!black}{red!5}{\parbox{0.92\textwidth}{
\textbf{#1}\\[4pt]#2
}}
\end{center}
}

\title{No-GPS Rough-Terrain Route Repeat\\
\large Field Evidence, System Design, and Validation Plan}
\author{}
\date{July 16, 2026}

\begin{document}
\maketitle
\tableofcontents
\newpage

\warningbox{Current Status}{
The no-GPS route-repeat pipeline is implemented and can send live rover
commands. It has \textbf{not} demonstrated a reliable autonomous repeat of the
taught route on the physical terrain. This is an experimental prototype, not a
cleared final-competition autonomy system.}

\section{Why This Is A Separate Branch}

The earlier outdoor system assumed a GPS mission: obtain checkpoint coordinates,
optionally expand the route with OpenStreetMap, execute a local controller, and
guard it with safety and recovery logic. The current session does not satisfy
that assumption. Telemetry reported placeholder GPS coordinates
\texttt{(1000, 1000)}, so a GPS controller cannot provide a meaningful global
navigation target.

The problem therefore changed from GPS mission execution to visual
teach-and-repeat:

\begin{center}
manual traversal $\rightarrow$ visual route memory $\rightarrow$ online visual
localization $\rightarrow$ local route-follow command
\end{center}

This is not generic exploration and it is not direct flag-seeking autonomy. It
is a route-specific repeated-navigation problem for a differential-drive rover
on rocky, uneven terrain.

\section{The System Contract}

\subsection{Sensors And Assumptions}

The branch deliberately has a narrow sensor contract:

\begin{itemize}[nosep]
  \item The live route-repeat runtime uses the \textbf{front camera only} through
  \texttt{/v2/front}.
  \item Rear camera imagery is not fused into localization or control.
  \item GPS is treated as unavailable for route following.
  \item IMU orientation, gyro rate, and RPM are support signals for tilt,
  turn-rate, and motion/stall reasoning; they are not a full pose estimator.
  \item The route is manually taught before the repeat attempt.
\end{itemize}

The front-only decision is intentional. Mixing an unvalidated rear stream into
the active localization input would make the visual reference problem less
clear, not more robust.

\subsection{Pipeline}

The implemented route pipeline is:

\begin{enumerate}[nosep]
  \item \texttt{scripts/record\_sdk\_session.py} records a manual traversal into
  an H5 bag. Front capture is the default.
  \item \texttt{tools/extract\_h5\_dataset.py} turns the H5 bag into ordered
  front images and metadata.
  \item \texttt{baseline.py} builds the descriptor database, place graph, and
  directed navigation graph.
  \item \texttt{scripts/prepare\_manual\_route.py} packages those stages into a
  reusable route directory.
  \item \texttt{scripts/run\_prepared\_route.py} launches the live runtime from
  that route package.
  \item \texttt{live\_indoor\_runtime.py} performs online localization, graph
  planning, control, recovery, and optional command sending.
\end{enumerate}

The runtime name is historical. In this branch it acts as the live no-GPS
route-repeat runner, not as the original indoor competition-only system.

\section{Route Data And What It Can Prove}

\begin{longtable}{p{0.27\textwidth}p{0.25\textwidth}p{0.38\textwidth}}
\toprule
\textbf{Artifact} & \textbf{Role} & \textbf{Current interpretation} \\
\midrule
\texttt{run\_01\_1521.h5} & Main manual teach bag & Strongest route-coverage candidate; not proof of autonomous repeat. \\
\texttt{smoke\_run01\_c} & 126-frame route package, target step 125 & Main reference and comparison package; not a deployment route. \\
\texttt{run\_01\_1521\_pp\_v3} & 15-frame post-processing experiment & Useful evidence that motion-aware selection is better; too sparse for primary use. \\
\texttt{run\_02\_1537.h5} & Alternate bag & Weaker due to a dark segment and poor tilted tail. \\
\texttt{run\_03\_1551.h5} & Alternate bag & Usable comparison only; weaker coverage than run 1. \\
\texttt{run\_05\_clean\_teach.h5} & Later recording & Saved but not yet post-processed or accepted as a route reference. \\
\bottomrule
\end{longtable}

\subsection{What The Bags Do Not Contain}

The manually recorded bags have an empty \texttt{controls} stream. They preserve
video, telemetry, IMU, magnetometer, and RPM evidence, but they do not preserve
the human driver's original linear and angular commands. Therefore, exact teleop
replay or direct command imitation cannot be reconstructed from the current
bags. Visual route following and command imitation are different problems.

\section{Localization, Planning, And Control}

\subsection{Visual Localization And Route Progression}

Each front frame is encoded against the taught route's descriptor database.
Temporal logic and a route corridor constrain the retrieved step, while the graph
planner selects a forward subgoal. This produces a controller input that includes
the current step, subgoal step, confidence, visual context, tilt, turn-rate, and
RPM motion information.

The conceptual loop is:

\begin{lstlisting}
front frame -> visual localization -> temporal/corridor check -> graph subgoal
            -> local controller -> terrain/recovery guard -> rover command
\end{lstlisting}

The graph step is a useful progress signal, but it is not physical ground truth.
It can remain unchanged while the rover moves if the live image does not match
the reference route well enough.

\subsection{Adaptive Pursuit And Runtime Recovery}

Rough-terrain prepared-route runs select
\texttt{src/adaptive\_pursuit\_controller.py}. The controller is intended to
move away from a brittle align-then-drive policy by combining forward motion with
continuous curvature-like steering. Its speed and turn behavior are conditioned
by:

\begin{itemize}[nosep]
  \item localization confidence,
  \item current and subgoal route steps,
  \item heading/turn-rate support signals,
  \item tilt,
  \item RPM motion hints,
  \item and no-progress state.
\end{itemize}

The surrounding runtime state machine adds startup probing, no-progress
relocalization, tilt-aware speed scheduling, stall handling, and jump rejection.
This is important: the observed robot behavior is the controller \emph{plus}
the recovery state machine, not the controller alone.

\subsection{Route Heading}

Route-heading guidance was tested, but compass/route-heading behavior was not
reliable in the field. In rough-terrain prepared-route mode,
\texttt{scripts/run\_prepared\_route.py} now defaults to
\texttt{--no-route-heading}. \texttt{--use-route-heading} is an explicit
opt-in for a route where it has been independently validated.

\section{What The Field Trials Actually Showed}

The field trials established several facts.

\begin{itemize}[nosep]
  \item The SDK connection, camera path, model loading, localization loop, and
  command-sending path all executed.
  \item The rover did physically move under \texttt{--send-control}.
  \item Localization repeatedly remained around startup steps 0, 1, or 8, often
  with moderate confidence in the approximate range 0.49--0.59.
  \item The runtime repeatedly entered \texttt{relocalize\_scan} and
  \texttt{relocalize\_probe\_forward} cycles without sustained graph-step
  progression.
  \item At roughly 14--22 degrees of tilt, cautious crawl commands often lacked
  enough authority to cross rocks, while stronger probes still did not create a
  durable reference-route match.
  \item A tightly constrained startup corridor prevented unsafe global jumps,
  but it could not create a valid localization match when the live start did not
  visually agree with the taught reference.
\end{itemize}

The correct conclusion is not that the rover ``did nothing.'' It moved. The
failure was that physical motion, visual localization, and graph progression did
not remain mutually consistent enough to prove route following.

\section{Implemented Does Not Mean Proven}

\begin{longtable}{p{0.36\textwidth}p{0.19\textwidth}p{0.35\textwidth}}
\toprule
\textbf{Capability} & \textbf{State} & \textbf{Meaning} \\
\midrule
Manual recording and route preparation & Implemented & A teach bag can become a visual route package. \\
Front-camera live localization & Implemented & The runtime receives and queries live front frames. \\
Adaptive-pursuit route command path & Implemented & Commands are computed and can be sent. \\
Terrain-aware recovery & Implemented / heuristic & Recovery behavior exists but is not field-proven to resolve localization failure. \\
Reliable autonomous taught-route repeat & Unproven & No controlled physical test has shown repeatable sustained progression. \\
Final-competition autonomous deployment & Not justified & The required validation evidence does not yet exist. \\
\bottomrule
\end{longtable}

\section{The Only Valid Next Milestone}

The next goal is not ``make the rover move more'' and not ``add another model.''
It is to prove one short, repeatable hardware-in-the-loop segment on the exact
physical route:

\begin{enumerate}[nosep]
  \item Place the rover where the physical scene visibly matches the taught
  route start.
  \item Run the route follower in dry-run and confirm that localization stays in
  the expected route corridor.
  \item Enable supervised control only after the dry-run output is coherent.
  \item Require meaningful forward graph-step progression without repeated
  scan/probe cycling.
  \item Repeat the same short segment successfully at least once.
\end{enumerate}

Failure to meet any one of these conditions means the branch remains an
experimental evaluation system. Moving physically is not enough; the movement
must be consistent with the taught route and reproducible.

\section{What Carries Forward From Older Work}

The no-GPS branch is distinct from the older GPS mission and marathon systems,
but it inherits their most important systems lessons:

\begin{itemize}[nosep]
  \item transition stability matters more than adding layers on paper;
  \item a weak perception signal must not be promoted into false authority;
  \item a recovery behavior can make a system look active while preventing real
  progress;
  \item physical stability, traction, and tilt are part of the control problem;
  \item field evidence outranks an attractive architecture diagram.
\end{itemize}

\section{Source Of Truth}

For exact implementation details, trust active code. For maturity and field
claims, trust the recorded test evidence first. The concise current status is in
\texttt{obsidian\_vault/00 Home/Current No-GPS - Read This First.md}; detailed
evidence is in
\texttt{obsidian\_vault/01 Source of Truth/No-GPS Field Trial - Findings.md}.

\end{document}
